\section{Principal divisors}
\label{section-principal-divisors}

\noindent
The following definition is the analogue of
Divisors, Definition \ref{divisors-definition-principal-divisor}
in our current setup.

\begin{definition}
\label{definition-principal-divisor}
Let $(S, \delta)$ be as in Situation \ref{situation-setup}.
Let $X$ be locally of finite type over $S$. Assume $X$ is
integral with $\dim_\delta(X) = n$.
Let $f \in R(X)^*$. The {\it principal divisor
associated to $f$} is the $(n - 1)$-cycle
$$
\text{div}(f) = \text{div}_X(f) = \sum \text{ord}_Z(f) [Z]
$$
defined in Divisors, Definition \ref{divisors-definition-principal-divisor}.
This makes sense because prime divisors have $\delta$-dimension $n - 1$ by
Lemma \ref{lemma-divisor-delta-dimension}.
\end{definition}

\noindent
In the situation of the definition for $f, g \in R(X)^*$ we have
$$
\text{div}_X(fg) = \text{div}_X(f) + \text{div}_X(g)
$$
in $Z_{n - 1}(X)$. See Divisors, Lemma \ref{divisors-lemma-div-additive}.
The following lemma will be superseded by the more general
Lemma \ref{lemma-flat-pullback-rational-equivalence}.

\begin{lemma}
\label{lemma-flat-pullback-principal-divisor}
Let $(S, \delta)$ be as in Situation \ref{situation-setup}.
Let $X$, $Y$ be locally of finite type over $S$. Assume $X$, $Y$
are integral and $n = \dim_\delta(Y)$.
Let $f : X \to Y$ be a flat morphism of relative dimension $r$.
Let $g \in R(Y)^*$. Then
$$
f^*(\text{div}_Y(g)) = \text{div}_X(g)
$$
in $Z_{n + r - 1}(X)$.
\end{lemma}

\begin{proof}
Note that since $f$ is flat it is dominant so that
$f$ induces an embedding $R(Y) \subset R(X)$, and hence
we may think of $g$ as an element of $R(X)^*$.
Let $Z \subset X$ be an integral closed subscheme of
$\delta$-dimension $n + r - 1$. Let $\xi \in Z$
be its generic point. If $\dim_\delta(f(Z)) > n - 1$,
then we see that the coefficient of $[Z]$ in the left and
right hand side of the equation is zero.
Hence we may assume that $Z' = \overline{f(Z)}$ is an
integral closed subscheme of $Y$ of $\delta$-dimension $n - 1$.
Let $\xi' = f(\xi)$. It is the generic point of $Z'$.
Set $A = \mathcal{O}_{Y, \xi'}$, $B = \mathcal{O}_{X, \xi}$.
The ring map $A \to B$ is a flat local homomorphism of
Noetherian local domains of dimension $1$.
We have $g$ in the fraction field of $A$. What we have to show is that
$$
\text{ord}_A(g) \text{length}_B(B/\mathfrak m_AB)
=
\text{ord}_B(g).
$$
This follows from Algebra, Lemma \ref{algebra-lemma-pullback-module}
(details omitted).
\end{proof}











